\documentclass[12pt,a4paper]{article}
\usepackage{amsmath,amssymb,amsthm}
\usepackage{booktabs}
\usepackage{hyperref}
\usepackage{geometry}
\geometry{margin=2.5cm}
\usepackage{enumitem}

\newtheorem{theorem}{Theorem}[section]
\newtheorem{lemma}[theorem]{Lemma}
\newtheorem{corollary}[theorem]{Corollary}
\newtheorem{proposition}[theorem]{Proposition}
\theoremstyle{definition}
\newtheorem{definition}[theorem]{Definition}
\theoremstyle{remark}
\newtheorem{remark}[theorem]{Remark}

\title{Planck's Constant from Recognition Science:\\
$\hbar = \varphi^{-5}$ in RS-Native Units}

\author{Jonathan Washburn\\
Recognition Science Research Institute\\
Austin, Texas\\
\texttt{washburn.jonathan@gmail.com}}

\date{March 2026}

\begin{document}
\maketitle

\begin{abstract}
We derive Planck's reduced constant $\hbar$ from Recognition
Science (RS) with zero free parameters.  In RS-native units
($\ell_0 = \tau_0 = 1$), the action quantum is:
$\hbar = E_{\mathrm{coh}} \cdot \tau_0 = \varphi^{-5}$,
where $E_{\mathrm{coh}} = \varphi^{-5}$ is the coherence energy
(the minimum energy for a recognition event) and $\tau_0 = 1$ is
the tick duration.  The exponent $5 = 2^D - D$ counts the compact
modes of the 8-tick epoch ($D = 3$).  We prove $\hbar > 0$,
$0.088 < \hbar < 0.093$ (in RS-native units), and the
quantum-gravitational duality $G \cdot \hbar = 1$, where
$G = \varphi^5$ is Newton's constant.  The SI value
$\hbar = 1.055 \times 10^{-34}\;\mathrm{J\,s}$ is recovered
through an external calibration seam (identification of
$\tau_0$ with the Planck time $t_P$).  $\hbar$ is not an
independent constant but a derived quantity---the minimum action
for a single recognition event.
\end{abstract}

%% ============================================================
\section{Introduction}
\label{sec:intro}
%% ============================================================

Planck's constant $\hbar = h/(2\pi)$ is the quantum of action,
setting the scale at which quantum effects become important.  In
conventional physics it is a measured constant; no theory derives
its value from more fundamental principles.

Recognition Science~\cite{RS_Axioms} derives $\hbar$ as a
$\varphi$-power.  The derivation has two ingredients: the
coherence energy $E_{\mathrm{coh}}$ (a property of the
$\varphi$-ladder) and the tick duration $\tau_0$ (a property of
the 8-tick epoch).  Their product is the minimum action for a
recognition event.

%% ============================================================
\section{Recognition Science Framework}
\label{sec:framework}
%% ============================================================

Recognition Science derives all physics from a single primitive,
the \emph{recognition cost functional} $J(x)$, uniquely determined
by three axioms.

\begin{definition}[RS Cost Functional]
For $x > 0$: $J(x) = \tfrac{1}{2}(x + x^{-1}) - 1$.
\end{definition}

\noindent\textbf{Axiom A1 (Normalization):} $J(1) = 0$.\\
\textbf{Axiom A2 (Recognition Composition Law, RCL):}
$J(xy) + J(x/y) = 2J(x)J(y) + 2J(x) + 2J(y)$ for all $x, y > 0$.\\
\textbf{Axiom A3 (Calibration):} $J''_{\log}(0) = 1$, where
$J_{\log}(t) := J(e^t)$.

\begin{theorem}[Cost Uniqueness]
\label{thm:unique}
The unique continuous solution to \textup{(A1)--(A3)} is
$J(x) = \tfrac{1}{2}(x + x^{-1}) - 1$.
\end{theorem}

\begin{proof}[Proof sketch]
Setting $f(t) = J_{\log}(t) + 1$ and rewriting \textup{(A2)} gives
the d'Alembert equation $f(s+t) + f(s-t) = 2f(s)f(t)$.  By the
Acz\'el classification theorem~\cite{Aczel1966}, the unique continuous
solution with $f(0) = 1$ and $f''(0) = 1$ is $f(t) = \cosh t$.
\end{proof}

\noindent The golden ratio $\varphi = (1+\sqrt{5})/2$ satisfies
$J(\varphi) = \varphi - 1$ and is forced by the $\varphi$-ladder
self-similarity.  Spatial dimension $D = 3$ is forced by the gap-45
synchronization $\operatorname{lcm}(8, 45) = 360$.  The coherence
energy $E_{\rm coh} = \varphi^{-5}$ and Newton's constant $G = \varphi^5$
sit at opposite rungs of the ladder, which is the RS resolution of the
quantum-gravity hierarchy.

%% ============================================================
\section{The Coherence Energy}
\label{sec:ecoh}
%% ============================================================

\begin{definition}[Coherence Energy]
\label{def:ecoh}
The coherence energy is the minimum energy quantum for a
recognition event:
\begin{equation}
  E_{\mathrm{coh}} = \varphi^{-5},
\end{equation}
where $\varphi = (1+\sqrt{5})/2$ is the golden ratio.
\end{definition}

\begin{proposition}[Origin of the Exponent]
\label{prop:exponent}
The exponent $5$ is determined by $D = 3$:
\begin{equation}
  5 = 2^D - D = 8 - 3.
\end{equation}
This counts the number of compact (non-spatial) modes of the
8-tick epoch.
\end{proposition}

\begin{proof}[Structural argument]
The 8-tick epoch has $2^D = 8$ binary configurations.  Of these,
$D = 3$ correspond to the spatial dimensions (the ``visible''
directions along which particles propagate).  The remaining
$2^D - D = 5$ are compact modes---internal phase-space directions
that are integrated over when projecting to 3+1 spacetime.  The
coherence energy is the $\varphi$-suppression from integrating
over these 5 compact modes:
$E_{\mathrm{coh}} = \varphi^{-(2^D - D)} = \varphi^{-5}$.

\emph{Status note.}  The identification of the 5 remaining
configurations as ``compact modes'' integrated over to produce
$\varphi$-suppression is a structural interpretation, not yet a
fully rigorous derivation from the 8-tick Hamiltonian cycle.
The exponent $5 = 2^D - D$ for $D = 3$ is the key claim; the
exact mechanism of $\varphi$-suppression per compact mode is an
open problem in RS.
\end{proof}

\begin{remark}[$E_{\mathrm{coh}}$ Bounds]
\label{thm:ecoh-bounds}
$E_{\mathrm{coh}} > 0$ (since $\varphi > 0$) and
$0.088 < E_{\mathrm{coh}} < 0.093$: from
$1.61 < \varphi < 1.62$, we get $10.7 < \varphi^5 < 11.4$,
hence $0.088 < \varphi^{-5} < 0.093$.
The exact value is $\varphi^{-5} = (7 - 3\sqrt{5})/2 \approx 0.09017$.
\end{remark}

%% ============================================================
\section{Derivation of $\hbar$}
\label{sec:derivation}
%% ============================================================

\begin{proposition}[Planck's Constant --- Structural Identification]
\label{thm:hbar}
In RS-native units ($\tau_0 = 1$), the action quantum is identified as:
\begin{equation}
  \boxed{\hbar = E_{\mathrm{coh}} \cdot \tau_0 = \varphi^{-5}
  \approx 0.0902.}
\end{equation}
\end{proposition}

\begin{proof}[Structural argument]
The RS identification is that the quantum of action equals the minimum
energy times the minimum time.  Both are forced by $D = 3$:
\begin{enumerate}[label=(\roman*),nosep]
  \item $E_{\mathrm{coh}} = \varphi^{-5}$: the minimum energy
  quantum (Definition~\ref{def:ecoh}).
  \item $\tau_0 = 1$: the minimum time quantum (one tick).
\end{enumerate}
By dimensional analysis, $[\text{action}] = [\text{energy}] \times [\text{time}]$,
so $\hbar = \varphi^{-5} \cdot 1 = \varphi^{-5}$.

\emph{Status note}: This is a structural identification ($\hbar$ is
identified with $E_\mathrm{coh} \cdot \tau_0$) rather than a
derivation of why the quantum of action should equal the product of
the minimum energy and minimum time.  The non-trivial content is
that $E_\mathrm{coh} = \varphi^{-5}$ (itself a structural argument,
Proposition~\ref{prop:exponent}).  A rigorous derivation connecting
the RS action quantum to $\hbar$ via the path integral or commutation
relations from first principles is an identified open problem.
\end{proof}

%% ============================================================
\section{Quantum--Gravitational Duality}
\label{sec:duality}
%% ============================================================

\begin{proposition}[Quantum-Gravitational Duality]
\label{thm:duality}
Under the RS identifications $G = \varphi^5$ and $\hbar = \varphi^{-5}$
(complementary rungs on the $\varphi$-ladder):
\begin{equation}
  \boxed{G \cdot \hbar = \varphi^5 \cdot \varphi^{-5} = 1.}
\end{equation}
\end{proposition}

\begin{proof}
$\varphi^{5} \cdot \varphi^{-5} = \varphi^{5+(-5)} = \varphi^0 = 1$.
The RS content is that $G$ and $\hbar$ sit at equal and opposite rungs
of the $\varphi$-ladder ($G$ at rung $+5$, $\hbar$ at rung $-5$),
making their product the ladder identity.  The identification
$G = \varphi^5$ (derived separately in the companion paper on Newton's
constant) is assumed here.
\end{proof}

\begin{remark}
The identity $G \cdot \hbar = 1$ is the RS resolution of the
hierarchy problem.  In conventional physics, the gravitational
and quantum scales differ by $\sim 10^{38}$ in SI units---a
``hierarchy'' that demands explanation.  In RS-native units,
they are exact reciprocals on the $\varphi$-ladder: $G$ is
rung $+5$, $\hbar$ is rung $-5$.  The apparent hierarchy in
SI units is an artifact of the calibration seam, not a physical
puzzle.
\end{remark}

\begin{remark}[$G$ convention]
An alternative RS derivation (recognition-wavelength
extremum) yields $G = \varphi^5/\pi$, giving
$G \cdot \hbar = 1/\pi$ instead of~$1$ and
$\ell_P = 1/\sqrt{\pi}$.
The duality ($G$ and~$\hbar$ at opposite rungs) is
preserved; only the product acquires a geometric
prefactor.  See the companion paper on Newton's constant.
\end{remark}

\begin{corollary}[Planck Mass]
In RS-native units ($c = 1$):
\begin{equation}
  M_{\mathrm{Pl}} = \sqrt{\frac{\hbar c}{G}}
  = \sqrt{\frac{\varphi^{-5}}{\varphi^5}}
  = \varphi^{-5} \approx 0.090.
\end{equation}
\end{corollary}

%% ============================================================
\section{Path Integral Regularization}
\label{sec:path}
%% ============================================================

\begin{proposition}[Natural UV Cutoff]
\label{prop:uv}
The path integral is naturally regularized by the discrete
ledger.
\end{proposition}

\begin{proof}[Argument]
In Feynman's path integral, amplitudes are:
$\mathcal{A} = \int \mathcal{D}[x]\, e^{iS[x]/\hbar}$.
In RS, each history is a finite sequence of recognition events,
each contributing action $\geq \hbar$.  The minimum action
$\hbar = E_{\mathrm{coh}} \cdot \tau_0$ means no process can
contribute less than one energy quantum in less than one tick.
The lattice spacing $\ell_0$ and tick duration $\tau_0$ provide
the natural UV cutoff:
\begin{equation}
  \Lambda_{\mathrm{UV}} = \frac{1}{\ell_0}
  = \frac{c}{\tau_0 \cdot c} = \frac{1}{\tau_0}
  \quad (\text{in } c = 1 \text{ units}).
\end{equation}
There is no need for an ad hoc regularization scheme; the
discreteness of the ledger provides it.
\end{proof}

%% ============================================================
\section{SI Calibration}
\label{sec:si}
%% ============================================================

The RS derivation yields $\hbar = \varphi^{-5}$ as an identity
in RS-native units.  The SI value requires an external
calibration seam: one measured quantity ($\tau_0$ in seconds)
anchors the conversion.

The natural calibration identifies $\tau_0$ with the Planck time:
\begin{equation}
  \tau_0 = t_P = \sqrt{\frac{\hbar G}{c^5}}
  \approx 5.39 \times 10^{-44}\;\mathrm{s}.
\end{equation}
In RS, the Planck time is the fundamental tick because
$G \cdot \hbar = 1$ in RS-native units forces $\ell_P = \ell_0 = 1$.
The coherence energy in SI then equals the Planck energy:
\begin{equation}
  E_{\mathrm{coh}}^{\mathrm{SI}} = \frac{\hbar}{t_P}
  = \sqrt{\frac{\hbar c^5}{G}}
  = E_P \approx 1.956 \times 10^{9}\;\mathrm{J},
\end{equation}
and the RS-native value $\varphi^{-5} \approx 0.090$ is the
coherence energy \emph{in units of the Planck energy}: the
coherence quantum is approximately $9\%$ of the Planck energy.
Restoring SI:
\begin{equation}
  \hbar = E_{\mathrm{coh}}^{\mathrm{SI}} \cdot t_P
  = E_P \cdot t_P
  \approx 1.956 \times 10^{9} \times 5.39 \times 10^{-44}
  \approx 1.055 \times 10^{-34}\;\mathrm{J\,s},
\end{equation}
matching the CODATA value~\cite{CODATA2024}
$\hbar = 1.054\,571\,817 \times 10^{-34}\;\mathrm{J\,s}$.

\begin{remark}[Why the RS-native value is not in eV]
The dimensionless RS-native quantity is
$\hbar_{\mathrm{RS}} = \varphi^{-5} \approx 0.090$; this is
a ratio in Planck units, not an energy.  In SI, the coherence
quantum $E_{\mathrm{coh}} = E_P \approx 1.96\times10^9\,\mathrm{J}$
(the Planck energy), which when multiplied by $t_P$ gives
$\hbar_{\mathrm{SI}}$.  The often-cited ``$E_{\mathrm{coh}}
\approx 90\,\mathrm{MeV}$'' refers to the hadronic calibration
anchor used in mass derivations, not to the Planck-unit
interpretation here.  The two calibrations address different
physical scales; the RS framework accommodates both through
different choices of the calibration seam.
\end{remark}

\begin{remark}[Calibration Honesty]
The SI numerical value of $\hbar$ is not predicted by RS alone
---it requires one external anchor ($\tau_0$ in seconds).  What
RS predicts without any external input is the
\emph{dimensionless} structure: $\hbar = \varphi^{-5}$,
$G = \varphi^5$, $G \cdot \hbar = 1$, and the exponent
$5 = 2^D - D$.  These are the falsifiable, zero-parameter
results.
\end{remark}

%% ============================================================
\section{Relation to Other Constants}
\label{sec:relations}
%% ============================================================

In RS-native units ($\ell_0 = \tau_0 = 1$), the fundamental
constants form a coherent $\varphi$-ladder:
\begin{center}
\begin{tabular}{@{}lll@{}}
\toprule
\textbf{Quantity} & \textbf{RS-native value} &
\textbf{$\varphi$-rung} \\
\midrule
Speed of light $c$ & $\ell_0/\tau_0 = 1$ & 0 \\
Action quantum $\hbar$ & $\varphi^{-5}$ & $-5$ \\
Coherence energy $E_{\mathrm{coh}}$ & $\varphi^{-5}$ & $-5$ \\
Thermal quantum $k_R$ & $\ln\varphi$ & (logarithmic) \\
Newton's constant $G$ & $\varphi^5$ & $+5$ \\
\bottomrule
\end{tabular}
\end{center}

\begin{proposition}[$c$, $\hbar$, $G$ Identities]
\label{prop:cGhbar}
$c$, $\hbar$, and $G$ are related by three identities:
\begin{enumerate}[label=(\roman*),nosep]
  \item $c = 1$ (one voxel per tick),
  \item $G \cdot \hbar = 1$ (quantum-gravitational duality),
  \item The Planck length
  $\ell_P = \sqrt{\hbar G / c^3} = 1$ in RS-native units.
\end{enumerate}
\end{proposition}

%% ============================================================
\section{Comparison}
\label{sec:comparison}
%% ============================================================

\begin{center}
\renewcommand{\arraystretch}{1.3}
\begin{tabular}{@{}lll@{}}
\toprule
& \textbf{Standard Model} & \textbf{RS} \\
\midrule
$\hbar$ & Free parameter (measured) & $\varphi^{-5}$ (derived) \\
$G$ & Free parameter (measured) & $\varphi^5$ (derived) \\
$G \cdot \hbar$ in RS-native units & $---$ & $= 1$ (exact) \\
Hierarchy problem & Unexplained & Reciprocal rungs \\
UV regularization & Ad hoc & Natural ($\ell_0, \tau_0$) \\
Free parameters & $\geq 2$ ($\hbar, G$) & 0 \\
\bottomrule
\end{tabular}
\end{center}

\begin{remark}[Hierarchy in SI]
In SI, $G \cdot \hbar \approx 7.0 \times 10^{-45}\,\mathrm{m}^3\,\mathrm{kg}^{-1}\,\mathrm{s}^{-1} \cdot \mathrm{J\,s}$
has units and is not directly comparable to the dimensionless
RS-native identity $G_{\rm RS} \cdot \hbar_{\rm RS} = 1$.  The
``hierarchy problem'' refers to the fact that in SI units, $G$
and $\hbar$ appear at vastly different scales only because SI
uses human-defined reference units.  In Planck units, the
hierarchy disappears: $G \cdot \hbar = (\hbar c / E_P^2)
\cdot (\hbar) = \hbar^2 c / E_P^2$.  The RS resolution is that
this dimensionless ratio equals 1 by the $\varphi^5 \cdot \varphi^{-5}
= 1$ identity.
\end{remark}

%% ============================================================
\section{Predictions and Falsifiers}
\label{sec:predictions}
%% ============================================================

\begin{enumerate}
  \item \textbf{$\hbar = \varphi^{-5}$}: Dimensionless identity
  in RS-native units.  The exponent $5 = 2^D - D$ is fixed by
  $D = 3$.

  \item \textbf{$G \cdot \hbar = 1$}: Zero-parameter prediction.
  Any deviation would falsify the $\varphi$-ladder structure.

  \item \textbf{No running of $\hbar$}: $\hbar$ is a fixed
  $\varphi$-power, not scale-dependent.

  \item \textbf{Minimum action}: No physical process can have
  action less than $\hbar$.  At the Planck scale, discreteness
  corrections of order $(\ell_0/\lambda)^2$ may become
  observable.

  \item \textbf{Falsifier}: An independent derivation of
  $\hbar$ from $D = 3$ geometry yielding an exponent
  $\neq 5 = 2^D - D$ would falsify the compact-mode
  interpretation.
\end{enumerate}

%% ============================================================
\section{Conclusion}
\label{sec:conclusion}
%% ============================================================

Planck's constant is not a free parameter but the minimum action
quantum: $\hbar = E_{\mathrm{coh}} \cdot \tau_0 = \varphi^{-5}$
in RS-native units.  The exponent $5 = 2^D - D$ counts the
compact modes of the 8-tick epoch.  Newton's constant
$G = \varphi^5$ is the exact reciprocal, giving the duality
$G \cdot \hbar = 1$.  The apparent hierarchy between quantum
and gravitational scales is an artifact of the SI calibration
seam; in RS-native units, they are mirror-image rungs on the
$\varphi$-ladder.

%% ============================================================
\begin{thebibliography}{99}

\bibitem{Aczel1966}
J.~Acz\'el, \emph{Lectures on Functional Equations and Their Applications},
Academic Press, New York (1966).

\bibitem{CODATA2024}
E.~Tiesinga, P.~J.~Mohr, D.~B.~Newell, and B.~N.~Taylor,
``CODATA Recommended Values of the Fundamental Physical
Constants: 2022,'' Rev.\ Mod.\ Phys.\ \textbf{96}, 025004
(2024).

\bibitem{Feynman1948}
R.~P.~Feynman, ``Space-Time Approach to Non-Relativistic
Quantum Mechanics,'' Rev.\ Mod.\ Phys.\ \textbf{20}, 367
(1948).

\bibitem{RS_Axioms}
J.~Washburn, ``The Algebra of Reality: A Recognition Science Derivation
of Physical Law,'' \emph{Axioms} \textbf{15}(2), 90 (2026).
\texttt{doi:10.3390/axioms15020090}

\end{thebibliography}

\end{document}
